\section{Devine le nombre}
\begin{UPSTIexercice}{Devine le nombre - Version While, avec n vies}
	Un nombre mystère est choisi au hasard entre 1 et 100. Le joueur a \texttt{n} vies pour deviner ce nombre.

	Tout d'abord, le programme affiche un menu avec 4 niveux de difficulté, qu'il choisira en entrant un chiffre entre 1 et 4 :
	\begin{enumerate}
		\item Facile (10 vies)
		\item Moyen (7 vies)
		\item Difficile (5 vies)
		\item Impossible (0 vies)
	\end{enumerate}
	Ensuite, le programme choisit un nombre au hasard entre 1 et 100, puis demande au joueur de deviner ce nombre.

	\begin{itemize}
		\item Après chaque proposition, le programme indique si le nombre mystère est plus grand ou plus petit que la proposition.
		\item Le jeu continue jusqu'à ce que le joueur trouve le nombre mystère ou qu'il n'ait plus de vies.
		\item Si le joueur trouve le nombre mystère, le programme affiche un message de félicitations.
		\item Si le joueur n'a plus de vies, le programme affiche un message de défaite et révèle le nombre mystère.
		\item Le programme doit utiliser une boucle \texttt{while} pour gérer les tentatives du joueur.
	\end{itemize}
	Le programme \texttt{devine\_while.c}contient le code suivant :
	\begin{lstlisting}[language=c]
#include <stdio.h>
#include <stdlib.h>
#include <time.h>
int main(void) {
    int vies;
    int nombre_mystere;
    int proposition;
    int niveau;
   // A COMPLETER : Demander le niveau de difficulté et initialiser le nombre de vies
   // L'utilisation d'un switch case parait appropriée ici.
    // Initialiser le générateur de nombres aléatoires
    srand(time(NULL));
    nombre_mystere = rand() % 100 + 1; // Nombre entre 1 et 100
   // A COMPLETER : Boucle principale du jeu
   return 0;
}
\end{lstlisting}
	\UPSTIetape Compléter le programme \texttt{devine\_while.c} pour qu'il fonctionne comme décrit.
	\UPSTIetape Tester le programme avec chaque niveau de difficulté.
	\UPSTIetape Que se passe-t-il si le joueur choisit le niveau "Impossible" ? Pourquoi ?
	\UPSTIetape Que se passe-t-il si le joueur choisi un niveau invalide (par exemple 5) ? Pourquoi ?
\end{UPSTIexercice}

\begin{UPSTIcorrectionP}{Devine le nombre - Version While, avec n vies}
	\begin{lstlisting}[language=c]
#include <stdio.h>
#include <stdlib.h>
#include <time.h>

int main(void) {
    int vies;
    int nombre_mystere;
    int proposition;
    int niveau;

    printf("Choisissez un niveau de difficulte :\n");
    printf("1. Facile (10 vies)\n");
    printf("2. Moyen (7 vies)\n");
    printf("3. Difficile (5 vies)\n");
    printf("4. Impossible (0 vies)\n");
    printf("Votre choix : ");
    scanf("%d", &niveau);

    switch (niveau) {
        case 1:
            vies = 10;
            break;
        case 2:
            vies = 7;
            break;
        case 3:
            vies = 5;
            break;
        case 4:
            vies = 0;
            break;
        default:
            printf("Niveau invalide.\n");
            return 1;
    }

    // Initialiser le générateur de nombres aléatoires
    srand(time(NULL));
    nombre_mystere = rand() % 100 + 1; // Nombre entre 1 et 100

    int trouve = 0;
    while (vies > 0 && !trouve) {
        printf("Proposez un nombre entre 1 et 100 (%d vies restantes) : ", vies);
        scanf("%d", &proposition);
        vies--;

        if (proposition == nombre_mystere) {
            trouve = 1;
        } else if (proposition < nombre_mystere) {
            printf("C'est plus grand !\n");
        } else {
            printf("C'est plus petit !\n");
        }
    }

    if (trouve) {
        printf("Bravo, vous avez trouve le nombre mystere !\n");
    } else {
        printf("Perdu ! Le nombre mystere etait %d.\n", nombre_mystere);
    }

    return 0;
}
	\end{lstlisting}
\end{UPSTIcorrectionP}

\begin{UPSTIexercice}{Devine le nombre - Version For, avec n vies}
	\UPSTIetape Faire le même exercice que précédemment, mais en utilisant une boucle \texttt{for} pour gérer les tentatives du joueur.
\end{UPSTIexercice}

\begin{UPSTIcorrectionP}{Devine le nombre - Version For, avec n vies}
	Seule la boucle principale change par rapport à la version \texttt{while} (le reste du programme -- choix du niveau, tirage du nombre mystère -- est identique) :
	\begin{lstlisting}[language=c]
    int trouve = 0;
    for (int i = 0; i < vies && !trouve; i++) {
        printf("Proposez un nombre entre 1 et 100 (%d vies restantes) : ", vies - i);
        scanf("%d", &proposition);

        if (proposition == nombre_mystere) {
            trouve = 1;
        } else if (proposition < nombre_mystere) {
            printf("C'est plus grand !\n");
        } else {
            printf("C'est plus petit !\n");
        }
    }
	\end{lstlisting}
\end{UPSTIcorrectionP}

\begin{UPSTIexercice}{Devine le nombre - Version Do...While, avec n vies}
	\UPSTIetape Faire le même exercice que précédemment, mais en utilisant une boucle \texttt{do...while} pour gérer les tentatives du joueur.
	\UPSTIetape Tester le programme avec n = 0. Que constatez-vous ?
	\UPSTIetape Ce type de boucle vous parait-il pertinent pour ce jeu ? Pourquoi ?
\end{UPSTIexercice}

\begin{UPSTIcorrectionP}{Devine le nombre - Version Do...While, avec n vies}
	Là encore, seule la boucle principale change :
	\begin{lstlisting}[language=c]
    int trouve = 0;
    do {
        printf("Proposez un nombre entre 1 et 100 (%d vies restantes) : ", vies);
        scanf("%d", &proposition);
        vies--;

        if (proposition == nombre_mystere) {
            trouve = 1;
        } else if (proposition < nombre_mystere) {
            printf("C'est plus grand !\n");
        } else {
            printf("C'est plus petit !\n");
        }
    } while (vies > 0 && !trouve);
	\end{lstlisting}
	Avec \texttt{n = 0} : contrairement aux versions \texttt{while} et \texttt{for}, le \texttt{do...while} exécute son corps \textbf{au moins une fois} avant de tester la condition. Le joueur est donc invité à proposer un nombre même avec 0 vie, alors que les deux autres versions annoncent la défaite immédiatement sans rien demander. Ce comportement est en général \textbf{non désiré} ici : le \texttt{do...while} n'est pas adapté dès que le nombre d'itérations peut être nul, ce qui est le cas pour ce jeu (niveau « Impossible »).
\end{UPSTIcorrectionP}

\begin{UPSTIexercice}{Vérifier le choix d'un menu}
	\UPSTIetape Quelle boucle est la plus adaptée pour vérifier que l'utilisateur a bien entré un choix valide dans un menu ? Justifier.
	\UPSTIetape Vérifier votre choix avec l'enseignant.
	\UPSTIetape Modifier vos programme précédents pour utiliser la boucle choisie.
\end{UPSTIexercice}

\begin{UPSTIcorrectionP}{Vérifier le choix d'un menu}
	La boucle \texttt{do...while} est la plus adaptée : on a besoin de lire une première fois la saisie de l'utilisateur \emph{avant} de pouvoir tester si elle est valide, puis de recommencer tant qu'elle ne l'est pas. Une boucle \texttt{while} obligerait à dupliquer la lecture (une fois avant la boucle, une fois dans la boucle) ou à utiliser une valeur \og sentinelle \fg{} artificielle pour entrer dans la boucle.
	\begin{lstlisting}[language=c]
int niveau;
do {
    printf("Choisissez un niveau de difficulte (1 a 4) : ");
    scanf("%d", &niveau);
} while (niveau < 1 || niveau > 4);
	\end{lstlisting}
\end{UPSTIcorrectionP}

\subsection*{Liste des checks disponibles}
\begin{lstlisting}[language=bash,style=console]
check50 IUT-GEII-Annecy/squelettes/branch-2025/tp3/devine_le_nombre/while
check50 IUT-GEII-Annecy/squelettes/branch-2025/tp3/devine_le_nombre/for
check50 IUT-GEII-Annecy/squelettes/branch-2025/tp3/devine_le_nombre/do_while
\end{lstlisting}
